\documentclass{article}
\usepackage[utf8]{inputenc}

\usepackage{natbib}

\usepackage{amsmath,amsfonts,amsthm}
\DeclareMathOperator{\spann}{span}
\DeclareMathOperator{\pca}{pca}
\DeclareMathOperator{\var}{var}
\DeclareMathOperator{\rank}{rank}
\DeclareMathOperator{\tr}{tr}
\DeclareMathOperator{\argmax}{argmax}

\newtheorem{theorem}{Theorem}
\newtheorem{lemma}{Lemma}
\newtheorem{definition}{Definition}


\title{
    \large Calculus
}
\author{Khanh Nguyen}
\date{May 2023}

\begin{document}
\maketitle

\emph{As a physics student in high school, CS undergraduate in NTU, I never have a proper/formal introduction to calculus. This tutorial is my best effort to formalize calculus for real numbers and served as a cheat sheet as well}

\section{Derivative}

Let $f \in \mathcal{C}^1(\mathbb{R}^m, \mathbb{R}^n)$. In this tutorial, to simplify the analysis, we assume $f$ being continuously differentiable.

We define the derivative $D: \mathcal{C}^1(\mathbb{R}^m, \mathbb{R}^n) \to \mathbb{R}^m \to \mathcal{L}(\mathbb{R}^m, \mathbb{R}^n)$ \footnote{
    \textbf{function currying notation} \\
    $A \to B$: set of all functions from $A$ to $B$ \\
    $A \to B \to C = A \to (B \to  C)$ \\
}
where

\begin{itemize}
    \item $\mathcal{C}^1(\mathbb{R}^m, \mathbb{R}^n)$: set of all continuously differentiable functions from $\mathbb{R}^m$ to $\mathbb{R}^n$
    \item $\mathcal{L}(\mathbb{R}^m, \mathbb{R}^n)$: set of all linear maps from $\mathbb{R}^m$ to $\mathbb{R}^n$
\end{itemize}

The derivative $D$ is an operator acting on the set of all continuously differentiable functions from $\mathbb{R}^m$ to $\mathbb{R}^n$ and produce a linear map on every points on $\mathbb{R}^m$

The derivative of $f$ at $x_0 \in \mathbb{R}^m$ is the \emph{closest} linear map \footnote{the existence and uniqueness of $Df(x_0)$ is a corrolary of Taylor's Theorem} to $f$ at $x_0$ 

\begin{equation}
    f(x) = f(x_0) + Df(x_0) (x - x_0) + o(x - x_0)
\end{equation}

where $o(x - x_0) \subset \mathbb{R}^m \to \mathbb{R}^n$ is the set of functions strictly smaller than $||x - x_0||$ when $x$ approaches $x_0$, i.e $\lim_{x \to x_0} \frac{||o(x - x_0)||}{||x - x_0||} = 0$

$Df(x)$ is a linear map from $\mathbb{R}^m$ to $\mathbb{R}^n$ which coincide with the Jacobian matrix $J_{f(x)} \in \mathbb{R}^{n \times m}$

\begin{equation}
J_{f(x)} = \begin{bmatrix}
\frac{\partial f_1}{\partial x_1} & ... & \frac{\partial f_1}{\partial x_m} \\
... & ... & ... \\
\frac{\partial f_n}{\partial x_1} & ... & \frac{\partial f_n}{\partial x_m} \\
\end{bmatrix}    
\end{equation}

\section{Derivative of element-wise function}

Given a scalar function $\sigma \in \mathcal{C}^1(\mathbb{R}, \mathbb{R})$ with derivate $D\sigma(x) \in \mathcal{L}(\mathbb{R}, \mathbb{R})$ which is a 1d linear function, i.e $y = ax$ for $a, x \in \mathbb{R}$

We can generalize this function to $n$-dimensional $\Bar{\sigma}(x) \in \mathcal{C}^1(\mathbb{R}^n, \mathbb{R}^n)$ such that if $y = \Bar{\sigma}(x)$ for $x \in \mathbb{R}^n$, then we have

\begin{equation}
    y_i = \sigma(x_i)    
\end{equation}

We can calculate the Jacobian matrix of this function

\begin{equation}
J = \begin{bmatrix}
\frac{\partial y_1}{\partial x_1} & ... & \frac{\partial y_1}{\partial x_n} \\
... & ... & ... \\
\frac{\partial y_n}{\partial x_1} & ... & \frac{\partial y_n}{\partial x_n} \\
\end{bmatrix}    
\end{equation}

where

\begin{equation}
    J_{i, j} = 
    \begin{cases}
    D\sigma(x_i) &\text{if $i = j$} \\
    0 &\text{if $i \neq j$} \\
    \end{cases}
\end{equation}

Jacobian matrix is diagonal

\begin{equation}
J = \begin{bmatrix}
D\sigma(x_0) & ... & 0 \\
... & ... & ... \\
0 & ... & D\sigma(x_n) \\
\end{bmatrix}    
\end{equation}

\section{Chain Rule}

Chain rule is used when we can to calculate the derivative of a function composition. Let's derive the chain rule.

Given $f \in \mathcal{C}^1(\mathbb{R}^n, \mathbb{R}^k)$, $g \in \mathcal{C}^1(\mathbb{R}^m, \mathbb{R}^n)$, and $f \circ g \in \mathcal{C}^1(\mathbb{R}^m, \mathbb{R}^k)$. We will calculate $D(f \circ g)$ given $Df$ and $Dg$

Let $x, x_0 \in \mathbb{R}^m$ and $y = g(x), y_0 = g(x_0) \in \mathbb{R}^n$

$Df(y_0)$ is the derivative of $f$ at $y_0$, we have

\begin{equation}
    f(y) = f(y_0) + Df(y_0)(y - y_0) + o(y - y_0)
\end{equation}

Rewrite

\begin{equation}
    (f \circ g)(x) = (f \circ g)(x_0) + Df(g(x_0))(g(x) - g(x_0)) + o(g(x) - g(x_0))
\end{equation}

$Dg(x_0)$ is the derivative of $f$ at $x_0$, we have

\begin{equation}
    g(x) = g(x_0) + Dg(x_0)(x - x_0) + o(x - x_0)
\end{equation}

Hence

\begin{align*}
    (f \circ g)(x)  &= (f \circ g)(x_0) + Df(g(x_0))(Dg(x_0)(x - x_0) + o(x - x_0)) + o(g(x) - g(x_0)) \\
                    &= (f \circ g)(x_0) + Df(g(x_0)) Dg(x_0)(x - x_0) + Df(g(x_0)) o(x - x_0) + o(g(x) - g(x_0)) \\
\end{align*}


We can prove $Df(g(x_0)) o(x - x_0) + o(g(x) - g(x_0)) = o(x - x_0)$ by assuming $||Df(g(x_0))|| < \infty$ and $||Dg(x_0)|| < \infty$ \\


$Df(g(x_0)) o(x - x_0)$ is small

\begin{align*}
    ||Df(g(x_0)) o(x - x_0)||  &\leq ||Df(g(x_0))|| ||o(x - x_0)|| &\text{(operator norm)} \\
\end{align*}

We have
\begin{align*}
    \lim_{x \to x_0} \frac{||Df(g(x_0)) o(x - x_0)||}{||x - x_0||}  &\leq ||Df(g(x_0))|| \lim_{x \to x_0} \frac{||o(x - x_0)||}{||x - x_0||} \\
                                                                    &= 0 \\
\end{align*}


$o(g(x) - g(x_0))$ is small

\begin{align*}
    \frac{||o(g(x) - g(x_0))||}{||x - x_0||}    &= \frac{||o(g(x) - g(x_0))||}{||g(x) - g(x_0)||} \frac{||g(x) - g(x_0)||}{||x - x_0||} \\
                                                &= \frac{||o(g(x) - g(x_0))||}{||g(x)- g(x_0)||} \frac{||Dg(x_0)(x - x_0) + o(x - x_0)||}{||x - x_0||} \\
                                                &\leq \frac{||o(g(x) - g(x_0))||}{||g(x)- g(x_0)||} \frac{||Dg(x_0)(x - x_0)|| + ||o(x - x_0)||}{||x - x_0||} &\text{(Cauchy-Schwarz)} \\
                                                &\leq \frac{||o(g(x) - g(x_0))||}{||g(x)- g(x_0)||} \frac{||Dg(x_0)|| ||(x - x_0)|| + ||o(x - x_0)||}{||x - x_0||} &\text{(operator norm)} \\
                                                &= \frac{||o(g(x) - g(x_0))||}{||g(x)- g(x_0)||} ( ||Dg(x_0)|| + \frac{||o(x - x_0)||}{||x - x_0||} )
\end{align*}

We have
\begin{align*}
    \lim_{x \to x_0} \frac{||o(g(x) - g(x_0))||}{||x - x_0||}  &\leq \lim_{x \to x_0} \frac{||o(g(x) - g(x_0))||}{||g(x)- g(x_0)||} ( ||Dg(x_0)|| + \frac{||o(x - x_0)||}{||x - x_0||} \\
                                                                    &= 0 \\
\end{align*}

Therefore,

\begin{equation}
    (f \circ g)(x) = (f \circ g)(x_0) + Df(g(x_0)) Dg(x_0)(x - x_0) + o(x - x_0) \\
\end{equation}

\begin{theorem} [Chain Rule]
Given two differentiable function $f \in \mathcal{C}^1(\mathbb{R}^n, \mathbb{R}^k)$, $g \in \mathcal{C}^1(\mathbb{R}^m, \mathbb{R}^n)$, the derivative of the composition $f \circ g$ is $D(f \circ g)(x) = Df(g(x)) Dg(x)$
\end{theorem}

\section{Change of variables}

Given a Lebesgue integral on a measurable subset $E \subset \mathbb{R}^n$ of positive measurable function $f: E \to [0, \infty)$ 

\begin{equation}
    \int_{E} f
\end{equation}

Let $G: \Omega \to E$ be a parameterization satisfying certain conditions \footnote{$\mathcal{C}^1$ diffeomorphism} where $\Omega$. Then

\begin{equation}
    \int_{E} f = \int_{\Omega} (f \circ G) |DG|
\end{equation}

where $|DG|$ is the determinant of derivative of $G$






\bibliographystyle{plain}
\bibliography{references}

\end{document}
